\documentclass[10pt,a4paper]{amsart}
\usepackage[margin=19mm]{geometry}
\usepackage{amsmath,amssymb,mathtools}
\usepackage[scr=rsfs,cal=euler]{mathalfa}
\usepackage{xfrac}
\usepackage[unicode,hidelinks,bookmarks=false]{hyperref}
\newcommand{\ie}{i.e.\ }
\newcommand{\cf}{cf.\ }
\newcommand{\resp}{respectively\ }
\newcommand{\Subsection}{\subsection}
\providecommand{\nobreakdash}{\nobreak\hbox{-}}
\setlength{\emergencystretch}{2em}
\multlinegap0pt
\usepackage{float}
\usepackage{amssymb}
\usepackage{mathtools}
\usepackage{braket}
\usepackage[shortlabels]{enumitem}
\usepackage{xcolor}
\usepackage{tikz}
\newtheorem{proposition}{Proposition}
\usepackage{cleveref}
\crefname{proposition}{Proposition}{Propositions}
\DeclareMathOperator{\ehr}{Ehr}
\DeclareMathOperator{\nint}{N_{\mathrm{int}}}
\title{Fano and Reflexive Polytopes from Feynman Integrals}
\author{Leonardo de la Cruz, Pavel P. Novichkov, Pierre Vanhove}
\date{Source: arXiv:2512.10518v2 (CC BY 4.0)}
\begin{document}
\maketitle

\noindent\emph{Source and licence.} Fano and Reflexive Polytopes from Feynman Integrals. Version arXiv:2512.10518v2 (CC BY 4.0), distributed by arXiv under the Creative Commons Attribution 4.0 International licence, http://creativecommons.org/licenses/by/4.0/, as recorded in the arXiv licence metadata for this exact version and shown on the abstract page https://arxiv.org/abs/2512.10518v2. Full licence notice: \texttt{licenses/arxiv-2512.10518-CC-BY-4.0.txt}.\par\smallskip

\noindent\textit{Editorial scope.} Counted unit: the article's proposition proposition (source0.tex lines 1812--1820). The article's complete original proof of it is delivered with it (source0.tex lines 1821--1832, one \texttt{proof} environment of 12 lines). No mathematics is written, repaired or re-derived: the statement and the proof are the article's own lines, in the article's own notation, and no retained line is substituted or re-flowed.  Retained uncounted as the source-local context the proof needs: lines 940--943.  Nothing was omitted from any retained span, so the package carries no omission record.  The retained lines cite 1 works, whose entries are transcribed verbatim from the article's own bibliography source in the e-print and delivered with short editorial labels recorded in the item's reference mapping.  No retained line contains a figure environment, an image-inclusion command or drawing code, so no image is delivered and the package declares no asset dependency.  Editorial formatting only, in addition: an AMS \texttt{amsart} preamble that restates the article's own declarations and macros used by the retained lines, namely the environments \texttt{proposition} on separate counters, together with the standard packages those lines need.  The retained environments print in this document as Proposition 1 in order of appearance, whereas the article prints the same items with the numbering of its own full document.  The complete native source of the article as distributed in the arXiv e-print for this exact version is bundled unchanged as \texttt{sources/190825/source0.tex}.
\begin{equation}
\label{eq:n_int_ehr}
 \ehr_P(-t) = (-1)^{\dim P} \nint(tP) \, ,
\end{equation}

\begin{proposition}[cf.~{\cite[Exercise 4.10]{Beck:2015}}]
	Let \(P\) be a \(d\)-dimensional polytope, \(d \geq 1\). Then the number of interior lattice points of its dilations satisfies
	\begin{equation}
		\begin{aligned}
			\nint((d+1)P) & \geq 1, \\
			\nint((d+2)P) & > 1.
		\end{aligned}
	\end{equation}
\end{proposition}

\begin{proof}
	This is a straightforward application of Ehrhart theory (see, e.g.,~\cite[Chapters 3 and 4]{Beck:2015}). The Ehrhart polynomial of \(P\) can be written in terms of the so-called \(h^*\)-vector as
	\begin{equation}
		\ehr_P(t) = \sum_{i=0}^d h_i^* \binom{t+d-i}{d},
		\label{eq:ehr_hstar}
	\end{equation}
	where \(h_0^* = 1\) and \(h_i^* \geq 0\). Combining \cref{eq:ehr_hstar} with the reciprocity formula~\eqref{eq:n_int_ehr}, we find
	\begin{equation}
		\nint(kP) = \sum_{i=0}^d h_i^* \binom{k+i-1}{d} \geq \binom{k-1}{d}. 
	\end{equation}
	Taking \(k=d+1\) and \(k=d+2\), we obtain the desired bound.
\end{proof}

\begin{thebibliography}{99}
\bibitem[Beck:2]{Beck:2015} M.~Beck and S.~Robins, \emph{Computing the Continuous Discretely: Integer-point Enumeration in Polyhedra}, Undergraduate Texts in Mathematics. Springer, New York, 2nd~ed., 2015, \href{https://doi.org/10.1007/978-1-4939-2969-3}{10.1007/978-1-4939-2969-3}.
\end{thebibliography}
\end{document}
